\documentclass[11pt]{article}
\usepackage{fontspec}
\setmainfont{Times New Roman}
\setsansfont{Arial}
\setmonofont{Consolas}
\usepackage{amsmath,amssymb}
\usepackage{geometry}
\usepackage{microtype}
\geometry{a4paper,margin=3cm}
\setlength{\parindent}{1.25em}
\setlength{\parskip}{0pt}
\emergencystretch=30em
\allowdisplaybreaks
\raggedbottom
\newcommand{\ideal}[1]{\mathfrak{#1}}
\newcommand{\eqzero}[1]{\equiv 0\,(#1)}

\begin{document}

% Unit: Paper 20. Direct Interslavic Latin translation from the German final-audited source slice.

\setcounter{footnote}{0}
\section*{20. Algebraičny kriterij absolutnoj nerazložimosti}
\begin{center}
Math. Ann. 85 (1922), str. 26--33
\end{center}

Polinom -- s koeficientami iz libovoljnogo, abstraktno opreděljenogo polja -- nazyvaje se \emph{absolutno nerazložimy}, ako ostavaje nerazložimy v algebraično zamknjenom polju, do ktorogo oblast koeficientov možno razširiti.\footnote{Že vsako polje, i to v suščnosti jednoznačno, možno razširiti do algebraično zamknjenogo polja, t. j. do takogo, v ktorom vsaky polinom jednej prěměnnoj razpadaje se na linearne faktory, pokazal E. Steinitz v svojej \emph{Algebraischen Theorie der Körper} (J. f. M. 137 (1910), str.~167); tym on dal racionalny ekvivalent fundamentalnogo teorema algebry.} V slědujučem pokažem, že absolutnu nerazložimost, v protivnosti k nerazložimosti vzhodno k zadannomu polju, možno algebraično shvatiti; točnejše važi teorem:

\emph{Vsakomu polinomu od $n\ge2$ prěměnnyh s neopreděljenymi koeficientami možno prirěditi cělu racionalnu funkciju s celočislenymi koeficientami od tyh koeficientov i daljnyh neopreděljenyh, formu reducibilnosti; tako, že za vsaky specialny polinom jednogo stepena nužno i dostatočno uslovje absolutnoj nerazložimosti jest dano neizčezanjem formy reducibilnosti. Drugymi slovami: za vsaku specialnu sistemu vrědnostij koeficientov, za ktoru forma reducibilnosti identično v neopreděljenyh izčezaje i ktora ne uslovjuje poniženje stepena, specialny polinom stavaje razložimy, i obratno.}

Ako vměsto polinomov smatrjajut se homogene formy v jednej prěměnnoj više, tada vměsto stepenovogo uslovja vystupaje prostějše uslovje, že koeficientna sistema ne izčezaje identično.

Kako neposrědno slědstvo teorema dobivaje se takože tvrđenje, najprěd postavjeno od \mbox{A.~Ostrowski}\footnote{\emph{Zur arithmetischen Theorie der algebraischen Größen}. Gött. Nachr. 1919, str.~279 (pomoćna lema, str.~296). Za slučaj poljev sostavljenyh iz obyčnyh čisel Ostrowski, koliko mi jest znano, takože došel do gornjogo glavnogo teorema; svoje izslědovanja o tom on skoro objavi. Že v mojem rasporedu dokaza glavny teorem važi za libovoljne, abstraktno opreděljene polja, opiraje se na tom, že forma reducibilnosti obrazovana za neopreděljene imaje čislove koeficienty nezavisne od specialno položeno v osnovu polja.}:

\emph{Vsaky absolutno nerazložimy polinom s algebraičnymi čislami kako koeficientami ostavaje nerazložimy modulo vsakogo prostogo ideala konečnogo algebraičnogo polja $\Omega$, ktoro sodrži oblast koeficientov, izključivši najviše konečno mnogo prostyh idealov iz $\Omega$. Osoblivo $\Omega$ može byti takože polje racionalnyh čisel.}

O daljnyh priměnjenjah k teoriji relativno cělyh funkcij\footnote{D. Hilbert, \emph{Mathematische Probleme}. Gött. Nachr. 1900, str.~253, problem 14.} naměrjam govoriti na drugom městu.

1. Najprěd pokažemo, že \emph{vsaky polinom od $n\ge2$ prěměnnyh s neopreděljenymi koeficientami jest absolutno nerazložimy}. Imenno, položimo:
\begin{equation}
\begin{aligned}
F(x)&=\sum A_{i_1\ldots i_n}x_1^{i_1}\cdots x_n^{i_n} &&(i_1+\cdots+i_n\le l),\\
G(x)&=\sum B_{j_1\ldots j_n}x_1^{j_1}\cdots x_n^{j_n} &&(j_1+\cdots+j_n\le m),\\
H(x)&=F(x)G(x)=\sum C_{k_1\ldots k_n}(A,B)x_1^{k_1}\cdots x_n^{k_n} &&(k_1+\cdots+k_n\le l+m),
\end{aligned}
\tag{1}
\end{equation}
kde $A$ i $B$ značat neopreděljene, a $C(A,B)$ stavajut celočislene bilinearne kombinacije od $A$ i $B$. Zato kvocienty
\[
D=C(A,B)/C_{0\ldots0}(A,B)
\]
sut racionalne funkcije kvocientov $A/A_{0\ldots0}$ i $B/B_{0\ldots0}$; ale čislo tyh funkcij stavaje večše než čislo jih argumentov; bo ono jest sootvětno
\[
\binom{l+m+n}{n}-1,
\qquad
\binom{l+n}{n}-1,
\qquad
\binom{m+n}{n}-1,
\]
i važi
\[
\binom{l+m+n}{n}+1>
\binom{l+n}{n}+\binom{m+n}{n}
\quad\text{za } n\ge2,\ l>0,\ m>0.
\]
Tako obstaja najmene \emph{jedna} algebraična relacija medžu $D(A,B)$ i slědovateljno takože medžu $C(A,B)$:
\begin{equation}
 \Phi(Z)\ne0\quad[\text{identično v }Z_{k_1\ldots k_n}],
 \qquad
 \Phi(C(A,B))=0\quad[\text{identično v }A,B],
\tag{2}
\end{equation}
kde $\Phi$ znači celočisleny polinom.

Polinom s neopreděljenymi $Z$ kako koeficientami:
\begin{equation}
 E(x)=\sum Z_{k_1\ldots k_n}x_1^{k_1}\cdots x_n^{k_n}
 \qquad(k_1+\cdots+k_n\le l+m)
\tag{3}
\end{equation}
jest zato za $n\ge2$ nužno absolutno nerazložimy.\footnote{Bo neopreděljene $Z$ po izrazě iz 2. ne sut nule od $\mathfrak P$.}

2. Cělost tyh polinomov $\Phi(Z)$, ktore pri substituciji $Z=C(A,B)$ identično v $A,B$ izčezajut, tvori \emph{ideal iz polinomov},\footnote{Take cělosti obyčno nazyvajut se moduly; faktično pak jih karakterizujuče svojstvo -- že one, pored $\Phi_1$ i $\Phi_2$, sodržat takože $\Phi_1-\Phi_2$, a pored $\Phi$ takože $\Psi\Phi$, kde $\Psi$ znači libovolny celočisleny polinom v $Z$ -- jest idealno svojstvo.} točnejše prosty ideal $\mathfrak P$; bo ako črez substituciju izčezaje produkt $\Phi_1\Phi_2$, izčezaje najmene jeden faktor. Poněže (2) jest izpolnjeno identično v $A,B$, ono ostavaje v sile za vsaku specialnu sistemu vrědnostij $A,B$, kde $A,B$ i slědovateljno takože $\Gamma=C(A,B)$ značat veličiny někojego polja. \emph{Koeficienty $\Gamma$ vsakogo polinoma razložimogo v razširitelnom polju zato zadovoljut uslovjam $\Phi(\Gamma)=0$, sut nule od $\mathfrak P$,} ili takože od konečno mnogih bazisnyh polinomov $\mathfrak P$; nyně ide o dokaz obrata.

K tomu treba primětiti, že vse relacije medžu $A,B,C$ ostavajut, ako polinomy (1) črez vvedenje novej prěměnnoj $y$ prěobrazujut se v homogene formy $F(x,y)$, $G(x,y)$, $H(x,y)$ stepenov $l,m,l+m$. Zato takože take koeficienty $\Gamma$ stavajut nulami $\mathfrak P$, pri ktoryh napr. $F(x,y)$\footnote{Polinomy i formy, ktore nastavajut, kogda neopreděljene v $F,G,\ldots,f,g,\ldots$ zaměnjajut se specialnymi sistemami vrědnostij, označam povsudu črtoju zgora: $\bar F,\bar G,\ldots,\bar f,\bar g,\ldots$.} stavaje stepenom od $y$; jim tedy črez prěhod k nehomogennym polinomam sootvětuje poniženje stepena od $\bar H(x)$, a ne razpadanje. Pokaže se, že izključivši to poniženje stepena obrat vsegda važi; osoblivo, že za \emph{homogene formy obrat važi bez izključkov, ako ne vse $\Gamma$ izčezajut; drugymi slovami, že vsaka nula od $\mathfrak P$ jest dana parametričnym predstavjenjem $\Gamma=C(A,B)$ s veličinami $A,B$ iz razširitelnogo polja}.\footnote{Že ``v obče'' nule od $\mathfrak P$ sut dane parametričnym predstavjenjem, slěduje iz svojstva prostogo ideala $\mathfrak P$ opreděljati nerazložimu algebraičnu mnogovidnost. Iz togo pak, kako znano, ne slěduje -- čto ja izvorno prěviděla i na čto \mbox{A.~Ostrowski} davno obratil moje vnimanje -- že parametrično predstavjenje ne može odpasti za žadnu specialnu sistemu vrědnostij. Tut dany dokaz opiraje se na elementarnějših pojmah.}

To dokazanje provedu direktnym obrazovanjem konečno mnogih funkcij $\Phi(Z)$; imennno črez Kroneckerovu substituciju
\[
 x_i=\xi^{d^{\,n-i}}
\]
prirědjam polinomam (1) polinomy od \emph{jednej} prěměnnoj; pokazujem, že tut nastavajut luky v eksponentah, i črez obrazovanje normy sootvětnyh koeficientov prihodim k funkcijam $\Phi(Z)$ i k iskannomu kriteriju.

3. Nehaj bude položeno:\footnote{\emph{Festschrift}, str.~11.}
\begin{equation}
\begin{aligned}
 d&=l+m+1,\\
 i&=i_1d^{n-1}+i_2d^{n-2}+\cdots+i_n,\quad
 j=j_1d^{n-1}+\cdots+j_n,\quad
 k=k_1d^{n-1}+\cdots+k_n.
\end{aligned}
\tag{4}
\end{equation}
Tada črez Kroneckerovu substituciju polinomam $F,G,H$ v (1) sootvětujut sootvětno polinomy
\begin{equation}
\begin{aligned}
 f(\xi)&=\sum A_{i_1\ldots i_n}\xi^i &&(i_1+\cdots+i_n\le l),\\
 g(\xi)&=\sum B_{j_1\ldots j_n}\xi^j &&(j_1+\cdots+j_n\le m),\\
 h(\xi)&=\sum C_{k_1\ldots k_n}(A,B)\xi^k &&(k_1+\cdots+k_n\le l+m).
\end{aligned}
\tag{5}
\end{equation}
To prirěđenje jest, kako znano, jedno-jednoznačno, poněže za vse v (5) vystupajuče indukovane eksponenty iz $k=0$ slěduje takože $k_1=0,\ldots,k_n=0$; i zato iz $i+j=i'+j'$ slěduje takože $i_1+j_1=i'_1+j'_1;\ldots;i_n+j_n=i'_n+j'_n$. Različne produkty $\xi^i\xi^j$ zato mogut dati ten sam eksponent $\xi^k$ samo togda, kogda i sootvětne produkty $x_1^{i_1}\cdots x_n^{i_n}x_1^{j_1}\cdots x_n^{j_n}$ vedut k tomu samomu eksponentu. Tako prihodit:
\begin{equation}
 h(\xi)=f(\xi)g(\xi),
\tag{6}
\end{equation}
i iz izpolnjenosti relacije (6), tako že v $f(\xi)$ i $g(\xi)$ ne vystupajut eksponenty različne od indukovanyh, nazad slěduje
\begin{equation}
 H(x)=F(x)G(x).
\tag{7}
\end{equation}
Pri tom \emph{v indukovanyh eksponentah v $f(\xi)$ i $g(\xi)$ v samoj stvari vystupajut luky}. Bo najvyšši eksponent od $f(\xi)$ stavaje $ld^{n-1}$, drugy najvyšši $(l-1)d^{n-1}+d^{n-2}$; razlika jest $d^{n-2}(d-1)>1$ zaradi $d=l+m+1\ge3$, $n\ge2$; i to samo važi za $g(\xi)$.

Relacije (6) i (7), važeče za neopreděljene $A,B$, ostavajut za vsaku specialnu sistemu vrědnostij $A,B$. Da by pri tom stepeni ostali, homogenizujem (6) i (7) črez prěměnne $\eta$ i $y$. \emph{Homogena forma $H(x,y)$ zato razpadaje se v algebraičnom razširitelnom polju togda i samo togda, kogda v njem sootvětna binarna forma $h(\xi,\eta)$ može se razložiti v dva faktory tako, že v faktorah ne vystupajut eksponenty različne od indukovanyh.}

4. Da by se provedlo obrazovanje normy pomenjeno na koncu 2., nehaj
\[
 a_1,b_1;\ldots; a_p,b_p;\quad a_{p+1},b_{p+1};\ldots; a_{p+q},b_{p+q}
\]
značat nove neopreděljene. Nehaj bude položeno
\begin{equation}
\begin{aligned}
 \varphi(\xi,\eta)&=(a_1\xi+b_1\eta)\cdots(a_p\xi+b_p\eta)
     =r_p\xi^p+r_{p-1}\xi^{p-1}\eta+\cdots+r_0\eta^p,\\
 \psi(\xi,\eta)&=(a_{p+1}\xi+b_{p+1}\eta)\cdots(a_{p+q}\xi+b_{p+q}\eta)
     =s_q\xi^q+s_{q-1}\xi^{q-1}\eta+\cdots+s_0\eta^q,\\
 \chi(\xi,\eta)&=\varphi(\xi,\eta)\psi(\xi,\eta)
     =t_{p+q}\xi^{p+q}+t_{p+q-1}\xi^{p+q-1}\eta+\cdots+t_0\eta^{p+q},
\end{aligned}
\tag{8}
\end{equation}
kde $r,s,t$ vsegda značat homogenizovane elementarne simetrične funkcije; t. j. te simetrične funkcije rędov $a,b$, linearnyh v vsakom jedinačnom rědu, iz ktoryh vse cěle simetrične funkcije, ktore sut v vsakom jedinačnom rědu homogene i jednakogo stepena, skladujut se cělo i s celočislenymi koeficientami -- i to samo jednym sposobom.

Nyně obrazujem s daljnymi neopreděljenymi $u_{\mu\nu}$ linearnu formu
\begin{equation}
 \zeta(u,a,b)=\sum u_{\mu\nu}r_\mu s_\nu,
\tag{9}
\end{equation}
gdě sumovanje razširjaje se po zadanyh indeksah $\mu$ i $\nu$. Produkt od $\zeta(u)$ so vsimi konjugatami, ktore nastavajut črez permutaciju rěda $a,b$ i sut različne od $\zeta(u)$ i medžu soboju -- norma $N(\zeta(u))$ od $\zeta(u)$ vzhodno k polju $t(a,b)$, ako $\zeta(u)$ smatra se kako funkcija polja od $r_\mu(a,b)$ i $s_\nu(a,b)$ -- stavaje potom celočislena simetrična funkcija vseh rędov $a,b$, v vsakom rědu homogena i jednakogo stepena, zato po gornjem jednoznačno opreděljena cěla celočislena funkcija od $t(a,b)$ i $u_{\mu\nu}$:
\begin{equation}
 N(\zeta(u,a,b))=T(t(a,b),u)\qquad[\text{identično v }a,b,u].
\tag{10}
\end{equation}
To samo razmysljenje pokazuje, že takože
\[
N(z-\zeta(u))=z^\delta+T_{\delta-1}(t,u)z^{\delta-1}+\cdots+T(t,u)
\qquad[\text{identično v }a,b,u,z]
\]
gdě vse $T$ značat cěle celočislene funkcije od $t$ i $u$. Obě strany toj identičnosti izčezajut za $z=\zeta(u)$; i to, zaradi algebraičnoj nezavisnosti elementarnyh simetričnyh funkcij $r(a,b)$, $s(a,b)$, identično v $r,s$, ako $t(a,b)$ po (8) položi se ravno celočislenej bilinearnoj kombinaciji $w(r,s)$ od $r$ i $s$, a $\zeta(u)$ smatra se kako funkcija od $r$ i $s$. Tako prihodit
\begin{equation}
\zeta(u)^\delta+T_{\delta-1}(w(r,s),u)\zeta(u)^{\delta-1}+\cdots+T(w(r,s),u)=0
\quad[\text{identično v }r,s,u].
\tag{11}
\end{equation}\footnote{(11), ako sumovanje v (9) razširjeno jest po vseh indeksah, predstavja Kroneckerovo razširenje Gaussovogo teorema, i to v suščnosti v Kroneckerovom rasporedu dokaza (\emph{Zur Theorie der Formen höherer Stufen}, Berliner Ber. 1883, Werke, 2, str.~417).}

Identičnosti (10) i (11) ostavajut za vse specialne sistemy vrědnostij $\alpha,\beta$ od $a,b$; sootvětno $\varrho,\sigma$ od $r,s$. Ako $\varrho,\sigma$ sut osoblivo izbrane tako, že vse produkty $\varrho_\mu\sigma_\nu$ izčezajut -- gde $\mu,\nu$ značat indekse vystupajuče v (9) -- tak že $\zeta(u)$ črez substituciju izčezaje, tada, ako položimo $w(\varrho,\sigma)=\tau$, po (11) prihodit:
\begin{equation}
 T(\tau,u)=0\qquad[\text{identično v }u].
\tag{12}
\end{equation}

Obratno, nehaj $\tau_0,\ldots,\tau_{p+q}$ sut take, ne vse izčezajuče sistemy vrědnostij, že (12) jest izpolnjeno. Poněže v algebraičnom razširitelnom polju obstaja razloženje
\begin{equation}
 \bar\chi(\xi,\eta)=\tau_{p+q}\xi^{p+q}+\cdots+\tau_0\eta^{p+q}
=(\alpha_1\xi+\beta_1\eta)\cdots(\alpha_{p+q}\xi+\beta_{p+q}\eta),
\tag{12'}
\end{equation}
kde v žadnom faktoru $\alpha$ i $\beta$ ne izčezajut jednočasno,\footnote{Toj (racionalny) ``fundamentalny teorem algebry'' jest eksistencijny teorem, na ktorom opiraje se bezizključkova validnost obrata izrečena na koncu 2.} ale něktore $\alpha$ ili $\beta$ mogut izčezati, stavaje $\tau=t(\alpha,\beta)$. Vstavjenje tyh vrědnostij v (10) pokazuje črez (12), že $\zeta(u,\alpha,\beta)$ ili jeden konjugatny faktor izčezaje identično v $u$. Ako zato, po vhodnoj numeraciji od $\alpha,\beta$, nyně položimo $r(\alpha,\beta)=\varrho$, $s(\alpha,\beta)=\sigma$, to znači, že $\bar\chi(\xi,\eta)$ dopuščaje najmene jedno razloženje:
\begin{equation}
\bar\chi(\xi,\eta)=\bar\varphi(\xi,\eta)\bar\psi(\xi,\eta)
=\sum \varrho_\mu\xi^\mu\eta^{p-\mu}\cdot\sum \sigma_\nu\xi^\nu\eta^{q-\nu},
\tag{13}
\end{equation}
tako, že vse produkty $\varrho_\mu\sigma_\nu$ vystupajuče v $\zeta(u)$ izčezajut, ale ne vse $\varrho$ ili $\sigma$.

5. Stepeni $p,q$ v (8) nehaj nyně budut ravne $ld^{n-1}$ i $md^{n-1}$, t. j. stepenam od $f(\xi)$ i $g(\xi)$ v (5). Indeksy $\mu,\nu$ razdělimo na $\mu^{(1)},\nu^{(1)},\mu^{(2)},\nu^{(2)}$; pri tom $\mu^{(1)}$ prohodi vse eksponenty, ktore chybajut v $f(\xi)$, $\nu^{(1)}$ vse eksponenty od $\xi$ v $\psi(\xi,\eta)$; i sootvětno $\nu^{(2)}$ vse eksponenty, ktore chybajut v $g(\xi)$, $\mu^{(2)}$ vse eksponenty od $\xi$ v $\varphi(\xi,\eta)$. Nadalje nehaj
\[
 e(\xi)=\sum Z_{k_1\ldots k_n}\xi^k\qquad(k_1+\cdots+k_n\le l+m)
\]
znači polinom od jednej prěměnnoj sootvětajuči polinomu (3); i nehaj
\[
 S(Z,u)
\]
znači celočisleny polinom v $Z$ i $u$, ktory nastavaje iz $T(t,u)$ -- jednoznačno opreděljene pravej strony (10), kde sumovanje v (9) razširjaje se po ukazanyh indeksah $\mu^{(1)},\nu^{(1)},\mu^{(2)},\nu^{(2)}$ -- tym, že $t$ zaměnjajut se koeficientami $Z$ i sootvětnymi nulami od $e(\xi)$.

Ako nyně $r,s$ zaměnjajut se koeficientami $A,B$ i sootvětnymi nulami od $f(\xi)$ i $g(\xi)$, togda črez tu substituciju $\zeta(u)$ izčezaje pri ukazanom sumovanji. Nadalje $t=w(r,s)$ prěhodi v koeficienty $C(A,B)$ i sootvětne nule od $h(\xi)$; zato $T(t,u)$ prěhodi v $S(C(A,B),u)$. Po identičnosti (11) zato prihodit
\begin{equation}
 S(C(A,B),u)=0\qquad[\text{identično v }A,B,u].
\tag{14}
\end{equation}
Poněže (14) ostavaje za vsaku specialnu sistemu vrědnostij, koeficienty $\Gamma=C(A,B)$ takyh specialnyh $H(x)$ ili občih $H(x,y)$, ktore v razširitelnom polju razpadajut se v dva faktory, zadovoljut uslovju
\begin{equation}
 S(\Gamma,u)=0\qquad[\text{identično v }u].
\tag{15}
\end{equation}

Ako obratno za, ne vse izčezajuče, koeficienty $\Gamma$ polinoma $\bar H(x)$, sootvětno formy $\bar H(x,y)$, uslovje (15) jest izpolnjeno, togda po gornjem to jest jednoznačno s $T(\tau,u)=0$, gde pod $\tau$ razumějut se vsi koeficienty od $\bar h(\xi,\eta)$. Po (13) zato obstaja v razširitelnom polju od $\Gamma$ razloženje:
\[
 \bar h(\xi,\eta)=\bar f(\xi,\eta)\cdot\bar g(\xi,\eta),
\]
gdě v $\bar f$ i $\bar g$ zaradi ukazanogo sumovanja ne vystupajut žadne od indukovanyh različne eksponenty; i slědovateljno obstaja relacija
\[
 \bar H(x,y)=\bar F(x,y)\cdot\bar G(x,y).
\]
Iz togo slěduje, ako za $y=1$ nijeden faktor ne stavaje nultego stepena, osoblivo ako $\Gamma$ ne uslovjujut poniženje stepena od $H(x)$, daljna relacija
\[
 \bar H(x)=\bar F(x)\cdot\bar G(x).
\]

\emph{Uslovje (15) jest zato nužno i dostatočno za to, že $\bar H(x,y)$, i pri izključenju poniženja stepena takože $\bar H(x)$, razpadaje se v razširitelnom polju v dva faktory.}

Iz togo dalje slěduje, že $S(Z,u)$ ne može izčezati identično v $Z$ i $u$, bo v 1. pokazano jest, že za $n\ge2$ polinomy s neopreděljenymi kako koeficientami, a zato takože sootvětne homogene formy v jednej prěměnnoj više, sut absolutno nerazložime. Tada, pod $U$ razuměvajuči stepenove produkty od $u$, prihodit:
\[
 S(Z,u)=\sum \Phi_\lambda(Z)U_\lambda,
\]
kde $\Phi_\lambda(Z)$ po (14) stavajut polinomami iz $\mathfrak P$. Tym jest pokazano, že $\mathfrak P$ ne imaje drugyh nul než ty dane parametričnym predstavjenjem $\Gamma=C(A,B)$, gde $A$ i $B$ prinadležat algebraičnomu razširitelnomu polju od $\Gamma$. Obrat izrečeny na koncu 2. jest tym dokazany.

6. \emph{Uslovje absolutnoj nerazložimosti} možno nyně neposrědnje formulovati. Zato treba samo stepen od $E(x)$ v (3) razložiti na vse možne načiny v dva pozitivne členy $l+m$, i k vsakomu razloženju obrazovati formu $S(Z,u)$. Produkt po tyh formah
\begin{equation}
 R(Z,u)=\prod S(Z,u)
\tag{16}
\end{equation}
tvori tada formu reducibilnosti od $E(x)$; bo po gornjemu vsakomu razpadanju od $E(x)$ sootvětno $E(x,y)$ sootvětuje izčezanje jednogo faktora od (16), i obratno. Tym jest dokazany glavny teorem postavjeny v uvodu i skupaj pokazano, kako forma reducibilnosti može se v samoj stvari postaviti v konečno mnogih krokah.

7. Iz eksistencije formy reducibilnosti (16) neposrědnje slěduje teorem Ostrowski pomenjeny v uvodu. Nehaj
\[
 \bar E(x)=\sum \Gamma_{k_1\ldots k_n}x_1^{k_1}\cdots x_n^{k_n}
 \qquad(k_1+\cdots+k_n\le v)
\]
bude absolutno nerazložimy polinom s cělymi algebraičnymi čislami $\Gamma$ kako koeficientami; nehaj $\Omega$ znači konečno algebraično čislovo polje, ktoro sodrži vse $\Gamma$, a $\mathfrak p$ prosty ideal iz $\Omega$.

Tada forma reducibilnosti $R(Z,u)$ obrazovana za točny stepen od $\bar E(x)$ pri substituciji $Z=\Gamma$ ostavaje različna od nuly; stavaje
\[
 R(\Gamma,u)=\sum P_\lambda U_\lambda\ne0\qquad[\text{identično v }u].
\]
Ideal $\mathfrak r=(P_1,P_2,\ldots)$ jest zato dělivy samo črez konečno mnogo prostyh idealov iz $\Omega$. Ako tedy $\mathfrak p$ jest različny od tyh konečno mnogih, togda $R(\Gamma,u)$ modulo $\mathfrak p$ ne izčezaje identično; i poněže klasy ostatkov modulo $\mathfrak p$ znovu tvore polje, slěduje, že $\bar E(x)$ modulo $\mathfrak p$ jest nerazložimy, točnejše absolutno nerazložimy. To samo važi za $\bar E(x,y)$, takže ne nastavaje ani poniženje stepena modulo $\mathfrak p$.

Ako koeficienty od $\bar E(x)$ sut algebraično drobne čisla, tada $\mathfrak p$ treba takože vzęti različny od konečno mnogih prostyh idealov, ktore vhodjat v obči imenovatelj $\Lambda$; modulo vseh ostalyh prostyh idealov možno zaměniti $E(x)$ črez $\Lambda\bar E(x)$, takže gorne prědpostavky sut izpolnjene. Tym jest teorem dokazany.

\begin{flushleft}
Erlangen, avgust 1921.
\end{flushleft}
\begin{center}
(Dostavljeno 28. 8. 1921.)
\end{center}

\clearpage

\end{document}
